
\documentclass[11pt]{article}

\usepackage[margin=1in]{geometry}
\usepackage{graphicx}
\usepackage{amsmath}
\usepackage{booktabs}
\usepackage{siunitx}

%%%%%%%%%%%%%%%%%%%%
% Personal packages
\usepackage{amssymb}
\usepackage{blindtext}
\usepackage{caption}

% Caption below table 
\usepackage{threeparttable}

% Long table and floating behavior handling
\usepackage{longtable}
% \usepackage{placeins}
\usepackage{xurl}
\usepackage{hyperref}

\usepackage[authoryear,longnamesfirst]{natbib}

%%%%%%%%%%%%%%%%%%%%

\title{
    Hydrogen-supply and feedstock sensitivity in pyrolysis-based sustainable aviation fuel production with South Korean context\\[1em]
    \Large\textbf{Supporting Information}
}
\author{Hyein Jung, Kyeongsu Kim, Jong Min Lee, Tae Hoon Oh}
\date{}

\begin{document}

\maketitle

\renewcommand{\thesection}{S\arabic{section}}
\counterwithin{table}{section}
\renewcommand{\thetable}{\thesection-\arabic{table}}

\section{Process Design Details}\label{sec:S0}

The pyrolysis process employs a circulating fluidized bed (CFB) reactor, selected for its scalability and proven performance in industrial applications by Ensyn \cite{ensyn_quebec} and Fortum \cite{fortum_biooil_joensuu}. 
The CFB reactor operates at 925~$^{\circ}$F, 21 psi with residence time in seconds to maximize bio-oil yield up to 75 wt\% while minimizing char formation.
Another key design parameter is the heat transfer regime, which utilizes sand as the fluidizing medium in this work. 
The sand enters the pyrolyzer with a temperature of 1,190~$^{\circ}$F and a sand-to-biomass ratio of 14.5 to ensure rapid heating rates and suppress secondary reactions. 

The pyrolysis product consists of three phases: crude bio-oil, solid char, and non-condensable gas (NCG). 
Crude bio-oil is a complex mixture of over 300 oxygenated compounds, modeled using representative groups from \cite{jones2013process}.
The model compounds used to represent each group are summarized in Table \ref{tab:model-compounds}. 
It is separated via a high-efficiency cyclone with 99.9\% capture rate, where heat-transferring sand is also separated.
Solid stream with char and sand is fed to the combustor, where the sand is heated up to 1189~$^{\circ}$F, ready to be sent back to the pyrolyzer.
NCG comprises CO, CO\textsubscript{2}, and light hydrocarbons as shown in Table \ref{tab:pyro-prod}.
A bi-level quenching system rapidly cools the vapor stream from 800~$^{\circ}$F to 113~$^{\circ}$F to stabilize reactive intermediates and remove NCG products from the bio-oil stream. 
The separated NCG components are compressed and fed back to the pyrolyzer as fluidizing gas. 

\begin{table}[htbp]
\centering
\caption{Composition of pyrolysis products, gas components, and bio-oil components.}
\begin{tabular}{lc||lc}
\toprule
\multicolumn{2}{c||}{\textbf{Pyrolysis Product}} & \multicolumn{2}{c}{\textbf{}} \\
\multicolumn{2}{c||}{(per dry biomass)} & \multicolumn{2}{c}{\textbf{Gas}} \\
\hline
\textbf{Components} & \textbf{wt\%} & \textbf{Components} & \textbf{wt\%}\\
\hline
Bio-oil & 64 & CO & 46 \\
Reaction water & 12 & CO\textsubscript{2} & 43 \\
Char \& ash & 12 & CH\textsubscript{4} & 6 \\
Gas & 12 & C\textsubscript{2}+ & 4.7 \\
& & H\textsubscript{2} & 0.3 \\
\midrule
\multicolumn{4}{c}{\textbf{Bio-oil Components}} \\
\hline
\textbf{Groups} & \textbf{wt\%} & \textbf{Groups} & \textbf{wt\%}\\
\hline
Acids & 5.57 & LMW Lignin A & 10.95 \\
Alcohols & 6.97 & LMW Lignin B & 2.21 \\
Ketones & 8.29 & Extractives & 2.99 \\
Aldehydes & 6.97 & HMW Lignin A & 9.15 \\
Guaiacols & 4.98 & HMW Lignin B & 1.99 \\
LMW Sugars & 5.97 & N-compounds & 0.07 \\
HMW Sugars & 33.86 & S-compounds & 0.025 \\
\bottomrule
\end{tabular}
\label{tab:pyro-prod}
\end{table}

Crude bio-oil requires upgrading to reduce its high oxygen content at 35~wt\% and acidity for fuel compatibility \cite{elliott2007historical}, achieved through a three-stage hydrotreating system \cite{elliott2016transportation}. First, a stabilizer operating at 140°C and 1200 psi with Ruthenium sulfide catalyst selectively hydrogenates reactive carbonyls and dienes, consuming minimal hydrogen (0.08 wt\%) to produce intermediates with lower acidity. The first-stage hydrotreater then processes moderate-reactivity oxygenates at 260°C and 2065 psi using Ru-based catalysts, requiring 7 wt\% hydrogen for temperature control via quenching and 0.18 wt\% for deoxygenation. Finally, sulfided molybdenum-based catalysts in the last stage (352°C, 2000 psi) saturate stubborn oxygenates and aromatics into cyclohexanes and iso-paraffins. Overall, the process converts 28.5 wt\% of dry biomass into upgraded oil with 12.2 wt\% hydrogen consumption. The composition of the upgraded oil was modeled  using user-defined components, based on experimental data reported in a previous study \cite{jones2013process}. Thermodynamic description of the user-defined components used in this study is in Table \ref{tab:app-user-define1}-\ref{tab:app-user-define2}.

After the sequential hydrotreating, the light column (T-290) separates NCG components smaller than C\textsubscript{5}.
In the design process, it separates pentane (light key component) from short-chain paraffins (heavy key components). 
The column employs a partial vapor condenser operating at 650 psi, which incorporates water entrainment to mitigate fouling and enhance separation performance. 
More detailed specification follows the description in Table \ref{tab:distil}.

\begin{table}[tbp]
\centering
\caption{Specifications for product distillation columns} \label{tab:distil}
\begin{threeparttable}
\begin{tabular}{l|c|c|c|c}
\toprule
\textbf{} & \textbf{T-290} & \textbf{T-310} & \textbf{T-320} & \textbf{T-330} \\
\midrule \midrule
Distillate & $<$ C\textsubscript{5} & Gasoline & SAF & Diesel \\
% \hline
% Light Key & Cyclo C6-HCK product & Cyclo C6 B & 1,2,3,4-Naphthalenes \\
% \hline
% Heavy Key & Cyclo C5 B & Phenols & Extractive intermediates \\
% \hline
No. of Stages & 30 & 30 & 20 & 20 \\
Feed stage & 8 & 15 & 10 & 10 \\
Reflux ratio\tnote{a} & 0.952 & 3.6 & 0.43 & 0.72 \\
Distillate/Feed\tnote{a} & 0.2 & 0.33 & 0.60 & 0.56 \\
Bottom Temp. & 367~$^{\circ}$F & 422~$^{\circ}$F & 512~$^{\circ}$F & 672~$^{\circ}$F \\
Condenser pressure & 65~psi & 45~psi & 20~psi & 15~psi \\
\bottomrule
\end{tabular}
\begin{tablenotes}
\footnotesize
\item[a] based on molar flow rate
\end{tablenotes}
\end{threeparttable}
\end{table}

    This study proposes sequential fractionation of hydrotreated bio-oil into three kinds of liquid fuels--gasoline, diesel, and SAF.
    % Upgraded bio-oil is fractionated into three kinds of liquid fuels - gasoline, diesel, and SAF. 
    The first column (T-310) for gasoline-range fuel operates at 45 psi, with cyclo C6-HCK products group as light key components and cyclo C5-B group as heavy components. 
    These hydrotreated compounds are defined in Table.\ref{tab:compounds-from-hydroprocess}.
    The last column (T-330) produces diesel-range fuel, separating 1,2,3,4-naphthalenes group from extractive intermediates.  
    The leftover long-chain hydrocarbons pass hydrocracker to produce short-chain products such as cyclo C6-HCK products group and return to the column T-310. 

%% Simulation: Pyrolysis_Jones_551_into_571_forQuality_TSOL2_retune.apw
\begin{table}[tb]
\centering
\caption{Comparison of Jet A and pyrolysis-based SAF product properties.}
\begin{threeparttable}
\begin{tabular}{l r r}
\toprule
\textbf{Property} & \textbf{Jet A} & \textbf{Pyrolysis SAF} \\
\midrule
\multicolumn{3}{l}{\textbf{Performance requirements}\tnote{a}} \\
% \midrule
Heat of combustion (MJ/kg) & $> 42.8 $ & 44.8 \\
Specific gravity  & 775--840  & 786 \\
Freezing point ($^\circ$C) & $< -40 $ & $-$56 \\
% Flash point, max ($^\circ$C) & $> 38 $ & 27.8 \\
Aromatic compounds (vol\%) & 8--25  & 10.2  \\
Sulfur content (wt\%) & $< 0.3 $ & NA \\
\midrule
\multicolumn{3}{l}{\textbf{Atomic composition comparison}\tnote{b}} \\
% \midrule
Carbon  (wt\%) & 86.4 & 83.8 \\
Hydrogen  (wt\%) & 13.5  & 13.6  \\
Oxygen  (wt\%) & 0.0044 & 0.026  \\
\bottomrule
\end{tabular}
\begin{tablenotes}
\footnotesize
\item[a] Ref. \cite{lahijani2022upgrading} 
\item[b] Ref. \cite{zhang2015production}
\end{tablenotes}
\end{threeparttable}
\label{tab:saf-spec}
\end{table}

    The second column (T-320) produces SAF within the C\textsubscript{8-16} hydrocarbon range, which is targeted to meet the performance specification of Jet A fuel \cite{Davidson2014}. 
    The evaluation of fuel quality was based on several key performance indicators: heat of combustion, flash point, specific gravity, and aromatic content \cite{lahijani2022upgrading}, as outlined in Table \ref{tab:saf-spec}.
    The heat of combustion was assessed conservatively using the lower heating value of the fuel components, while the flash point was determined by the lowest boiling point of the constituent compounds. 
    Specific gravity and volumetric fraction of aromatic compounds were evaluated via Aspen Plus simulations.
    The freezing point was estimated according to carbon class composition, employing correlation models proposed by \cite{shi2017quantitative}.
    Among these models, coefficients optimized by a genetic algorithm were selected, as they demonstrated the highest accuracy for freezing point estimation in catalytic fast pyrolysis-derived SAF \cite{watanasiri2023prediction}. 

    With the exception of the volumetric fraction of aromatic compounds, all evaluated properties were found to be within the desired range. 
    To address this, phenols were separated from the final product, which resulted in an aromatic content of 10.2 vol\%.
    Additionally, the atomic composition of carbon, hydrogen, and oxygen in the produced fuel was compared to that of conventional jet fuel \cite{zhang2015production}, confirming that the composition of T-320 product is comparable to standard aviation fuels.
    
    The off-gas collected from hydroprocesses, hydrotreating, and hydrocracking is rich in hydrogen gas.
    Hence, the leftover hydrogen gas is recovered via a pressure swing adsorption (PSA) unit to satisfy the high purity ($> 99$ wt\%) required for the hydroprocesses.
    Hydrogen recovery is assumed to be 85\%, referring to the industrial technology providing approximately 80\% recovery rate these days \cite{mahler_hydroswing}.
    The recovered hydrogen gas is recycled for the hydroprocesses with makeup feed from external sources. 
    
    The off-gas consists of 74.8 vol\% hydrogen gas, so the state-of-the-art gas turbine technology should be applied to resolve the high flame temperature and flame speed of hydrogen gas combustion.
    Exhaust gas recirculation (EGR) framework was applied to fulfill the maximum combustion temperature of 3002~$^{\circ}$F and prevent NOx emission \cite{dennis2024literature}.
    The gas turbine generates 65 MW with high-temperature waste heat, which is utilized in the distillation column reboilers, EGR preheating, and CCS absorber regeneration. 

Table \ref{tab:model-compounds} provides the information of bio-oil model compounds used in process modeling. 

\begin{table}[htbp]
\centering
\caption{Model compounds and their composition on a dry weight basis retrieved from \cite{jones2013process}}
\begin{tabular}{|p{4.5cm}|p{4.5cm}|l|r|}
\hline
\textbf{Compound group} & \textbf{Model compound} & \textbf{Formula} & \textbf{\% wt} \\
% \textbf{} & \textbf{} & \textbf{} & \textbf{(dry basis)} \\
\hline
\multicolumn{4}{|l|}{\textbf{\textit{Water soluble}}} \\
\hline
Acids & Crotonic acid & C\textsubscript{4}H\textsubscript{6}O\textsubscript{2} & 5.57 \\
\hline
Alcohols & 1,4-Benzenediol & C\textsubscript{6}H\textsubscript{6}O\textsubscript{2} & 6.97 \\
\hline
Ketones & Hydroxyacetone & C\textsubscript{3}H\textsubscript{6}O\textsubscript{2} & 8.29 \\
\hline
Aldehydes & 3-Methoxy-4-Hydroxybenzaldehyde & C\textsubscript{8}H\textsubscript{8}O\textsubscript{3} & 6.97 \\
\hline
Guaiacols & Isoeugenol & C\textsubscript{10}H\textsubscript{12}O\textsubscript{2} & 4.98 \\
\hline
Low MW sugars & Levoglucosan & C\textsubscript{6}H\textsubscript{10}O\textsubscript{5} & 5.97 \\
\hline
High MW sugars & Cellobiose & C\textsubscript{12}H\textsubscript{22}O\textsubscript{11} & 33.86 \\
\hline
\multicolumn{4}{|l|}{\textbf{\textit{Water insoluble}}} \\
\hline
Low MW lignin derived compound A & Dimethoxy stilbene & C\textsubscript{16}H\textsubscript{16}O\textsubscript{2} & 10.95 \\
\hline
Low MW lignin derived compound B & Dibenzofuran (representing diphenyl compounds) & C\textsubscript{12}H\textsubscript{8}O & 2.21 \\
\hline
Extractives & Dehydroabietic acid & C\textsubscript{20}H\textsubscript{30}O\textsubscript{2} & 2.99 \\
\hline
High MW lignin derived compound A & Oligomeric compounds with $\beta$-O-4 bond & C\textsubscript{20}H\textsubscript{26}O\textsubscript{8} & 9.15 \\
\hline
High MW lignin derived compound B & Phenylcoumaran compounds & C\textsubscript{21}H\textsubscript{26}O\textsubscript{8} & 1.99 \\
\hline
Nitrogen compounds & 2,4,6-trimethylpyridine & C\textsubscript{8}H\textsubscript{11}N & 0.070 \\
\hline
Sulfur compounds & Dibenzothiophene & C\textsubscript{12}H\textsubscript{8}S & 0.025 \\
\hline
\end{tabular}
\label{tab:model-compounds}
\end{table}

\begin{table}[htbp]
\centering
\caption{Boiling point and vapor pressure equation coefficients of user-defined compounds retrieved from \cite{jones2013process}.}
\begin{threeparttable}
% \scriptsize
\begin{tabular}{|l|l|l|p{4.5cm}|}
\hline
\textbf{Compound} & \textbf{Formula} & \textbf{Tbp, F} & \textbf{Coefficients for vapor pressure equation\tnote{1}} \\
\hline
Isoeugenol & C$_{10}$H$_{12}$O$_2$ & 507.2 & \begin{tabular}{l}
A = 50.75 \\
B = -9430.29 \\
C = -3.46 \\
D = -1.17E-6 \\
E = 2 \\
(Ref 2)
\end{tabular} \\
\hline
Levoglucosan & C$_6$H$_{10}$O$_5$ & 772.8 & \begin{tabular}{l}
A = 149.33 \\
B = -18972.6 \\
C = -16.826 \\
D = -4.57E-9 \\
E = 2 \\
(Ref 2)
\end{tabular} \\
\hline
Cellobiose & C$_{12}$H$_{22}$O$_{11}$ & 1234.2 & \begin{tabular}{l}
A = 53.01 \\
B = -19999.5 \\
C = -4.4294 \\
D = -1.5E-7 \\
E = 2
\end{tabular} \\
\hline
\begin{tabular}{l}
High MW lignin \\
derived compound A
\end{tabular} & C$_{20}$H$_{26}$O$_8$ & 1198.4 & \begin{tabular}{l}
A = 81.08 \\
B = -20000 \\
C = -7.036 \\
D = -2.279E-7 \\
E = 2
\end{tabular} \\
\hline
\begin{tabular}{l}
High MW lignin \\
derived compound B
\end{tabular} & C$_{21}$H$_{26}$O$_8$ & 1172.4 & \begin{tabular}{l}
A = 85.84 \\
B = -20000 \\
C = -7.72 \\
D = -3.84E-7 \\
E = 2
\end{tabular} \\
\hline
\end{tabular}
\begin{tablenotes}
\footnotesize
\item[1] Vp = exp[A + B/T + C ln(T) + D T$^E$], where Vp = vapor pressure (Pa), T = temperature (K)
\end{tablenotes}
\end{threeparttable}
\label{tab:app-user-define1}
\end{table}

\begin{table}[htbp]
\centering
\caption{Thermodynamic and critical properties of user-defined compounds retrieved from \cite{jones2013process}.}
% \scriptsize
\begin{tabular}{|l|l|p{1.1cm}|p{1.2cm}|p{1.1cm}|p{1.5cm}|p{1.5cm}|}
\hline
\textbf{Compound} & \textbf{Tc, F} & \textbf{Pc, psia} & \textbf{Vc, ft$^3$/ lbmol} & \textbf{Tm, F} & \textbf{$\Delta H_f$ (Btu/ lbmol)} & \textbf{$\Delta G_f$ (Btu/ lbmol)} \\
\hline
Isoeugenol  & 987.5 & 476.4 & 7.8 & 57.2 & -88516.7 & -16887.1 \\
\hline
Levoglucosan  & 1056 & 668.9 & 3.0 & 99.5 & -412740 & -255625 \\
\hline
Cellobiose  & 1281.5 & 340 & 6.8 & 462.2 & -958499 & -503968 \\
\hline
\begin{tabular}{l}
High MW \\ lignin derived \\ compound A
\end{tabular}  & 1482.2 & 241.9 & 15.9 & 218.3 & -498806 & -297590 \\
\hline
\begin{tabular}{l}
High MW \\ lignin derived \\ compound B
\end{tabular}  & 1468.2 & 243.7 & 16.0 & 252.3 & -496040 & -284159 \\
\hline
\end{tabular}
\label{tab:app-user-define2}
\end{table}

\FloatBarrier
\begin{longtable}{|p{0.4\textwidth}|p{0.5\textwidth}|}
\caption{Compound groups rising in hydroprocesses and their model compounds retrieved from \cite{jones2013process}.} \label{tab:compounds-from-hydroprocess}\\
\hline
\textbf{Groups} & \textbf{Model compounds} \\
\hline
\endfirsthead

\hline
\textbf{Generic names} & \textbf{Model compounds used to represent groups of similar compounds} \\
\hline
\endhead

1,2,3,4-Naphthalenes & 6-Hexyltetralin \\ \hline
Acids Intermediates & N-Butyric acid \\ \hline
Aromatic A & O-Xylene \\ \hline
Aromatic B & Benzene, 1-ethenyl-4-ethyl- \\ \hline
Ash & Ash \\ \hline
Bicyclos & 1,1-Bicyclohexyl \\ \hline
Carbon & Carbon \\ \hline
Carbonyl intermediates & 1,2-Propanediol \\ \hline
Extractive intermediates & HT Dehydroabietic acid \\ \hline
Cyclo C5 A & Ethyl-cyclopentane \\ \hline
Cyclo C5 B & 1-methyl-1-ethylcyclopentane \\ \hline
Cyclo C6 A & Cyclohexane \\ \hline
Cyclo C6 B & N-Butylcyclohexane \\ \hline
Cyclo C6-HCK product & Methyl-cyclohexane \\ \hline
Cyclo C7+ & 1,3-Dimethyladamantane \\ \hline
Diphenyl compounds & 1,2-Diphenylethane \\ \hline
Extractives & Dehydroabietic acid \\ \hline
High MW lignin compound intermediates & Indene \\ \hline
Indans & Indan \\ \hline
Indenes & 1H-Indene, 1,2,3-trimethyl- \\ \hline
Iso-paraffin A & 3-Methylhexane \\ \hline
Iso-paraffin B & 4-Methylnonane \\ \hline
Low MW lignin compound intermediates & Pinosylvin \\ \hline
Medium chain n-paraffins & N-Dodecane \\ \hline
Naphthalenes & 2,7-Dimethylnaphthalene \\ \hline
N-paraffin-HCK product A & N-Nonane \\ \hline
N-paraffin-HCK product B & N-Octane \\ \hline
Phenanthrenes & 4-Methylphenanthrene \\ \hline
Phenols & Thymol \\ \hline
Polynuclear aromatics & Naphthalene, 1-phenyl- \\ \hline
Pyrenes & Pyrene \\ \hline
Sand & Silicon Dioxide \\ \hline
Short chain n-paraffins & N-Hexane \\ \hline
Sugar intermediate A & Alpha-D-Glucose \\ \hline
Sugar intermediate B & Sorbitol \\ \hline
Sulfur agent & Dimethyl Sulfide \\
\hline
\end{longtable}

\newcommand{\COtwo}{CO\textsubscript{2}}
\newcommand{\COtwoe}{CO\textsubscript{2}e}
\newcommand{\Htwo}{H\textsubscript{2}}
\newcommand{\NHthree}{NH\textsubscript{3}}
\newcommand{\code}[1]{\texttt{#1}}

\section{Cost basis and financial parameters}\label{sec:S1}

\begin{itemize}
  \item Cost-year indices: \code{cost\_idx\_1Q2022 = 3230}, \code{cost\_idx\_Jones = 2250};
    Jones-era source costs escalated $\times(3230/2250)$ to 1Q-2022 USD. SMR uses its own index 2340.
  \item Financial: economic period 30 yr; IRR/discount 10\%; tax 35\%; salvage 0; straight-line
    depreciation; 8000 operating hours/yr ($\approx$91.3\% capacity factor).
  \item MJSP (annuity form):
    \begin{align*}
      \text{annuity} &= \frac{(1+\text{irr})^{-1}\,\bigl(1-(1+\text{irr})^{-n}\bigr)}{1-(1+\text{irr})^{-1}},\\
      \text{capital investment/yr} &= \frac{\text{CAPEX} - \text{annuity}\cdot\text{tax}\cdot\text{depreciation}}{\text{annuity}\,(1-\text{tax})},\\
      \text{MJSP} &= \frac{\text{min revenue} - \text{co-product sales}}{\text{annual SAF volume}}.
    \end{align*}
  \item SAF annual volume = stream 326SAF volumetric flow $\times$ 8000 h $\times$ 28.3168 L/cuft.
  \item Product/utility prices: 326SAF = 0 (the priced product); 316GAS (naphtha) 0.75757 \$/lb;
    336DIES (diesel) 0.6 \$/lb; electricity 19.05 \$/MMBtu; natural gas / fuel 14.38 \$/MMBtu.
\end{itemize}

%% ------------------------------------------------------------------
\section{Sub-unit cost interpolation}\label{sec:S2}

All five bolt-on sub-units use the same law
($r$ = stream/base stream): \textbf{CAPEX $\times r^{0.6}$} (six-tenths), \textbf{OPEX $\times r^{0.7}$},
stream and utility rates linear ($\times r$). Literature base cases are listed in Table~\ref{tab:S2-1}.
Hydrotreating catalyst charging: R-202/R-220 ruthenium LHSV 0.5,
price $60 \times 2.2$ \$/kg; R-225 molybdenum LHSV 0.22, price $15.5 \times 2.2$ \$/kg.

Its GHG emission is calculated from the simulated CCS cases and a constant ammonia-related emissions adjustment:
    \begin{equation}
    \mathrm{GHG}(a,c) =
    \mathrm{GHG}_{\mathrm{CCS}}(c)
    -a\left(\mathrm{GHG}_{\mathrm{grey}}-\mathrm{GHG}_{\mathrm{NH}_3}\right),
    \end{equation}
    where \(a\) denotes the green-ammonia substitution fraction and \(c\) denotes the CCS capture fraction. 

\begin{table}[htbp]
  \centering
  \caption{Literature base cases of the bolt-on sub-units.}
  \label{tab:S2-1}
  \resizebox{\textwidth}{!}{%
  \begin{tabular}{@{}llrrl@{}}
    \toprule
    Sub-unit & Reference & Base stream & Base CAPEX & Notes \\
    \midrule
    CCS (amine) & Rao \& Rubin 2002 & 862,853 lb/hr \COtwo & \$297.93 M & heat 3,315 MMBtu/hr; elec = 14\% of heat \\
    PSA \Htwo{} recovery & Byun 2020 & 700 m\textsuperscript{3}/hr & \$780,993 & OPEX \$5,078/yr; $\sim$68\% of OPEX = adsorbent/reactant \\
    R-510 hydrocracking & Jones 2013 & 9,600 lb/hr & \$5.27 M & driven by recycle stream 340 \\
    Off-gas reformer & Jones 2013 & (8846+995)/0.91 lb \Htwo & \$69.14 M (2011) & escalated $\times$3230/2250 \\
    SMR (+CCS) & Katebah 2022 & 45,922 & \$424 M / OPEX \$147 M & idx 2340; NG feed 156,136; \COtwo{} 55,107 \\
    \bottomrule
  \end{tabular}}
\end{table}


%% ------------------------------------------------------------------
\section{Scale-up and single-train limits}\label{sec:S3}

\begin{itemize}
  \item Scale set: [1, 2.5, 5, 7.5, 10, 20, 30]$\times$; main plant process 0.343~MMt/yr of lignocellulosic biomass.
  \item Pyrolyzer (R-110) single-train cap 700 MMBtu/hr = 10$\times$ base duty (70 MMBtu/hr at
    45.5 t/hr biomass); train count = ceil(base duty $\times$ scale / cap) $\rightarrow$ 1/1/1/1/1/2/3 across
    the scale set. Gas turbines, compressors, and expanders follow the same numbering-up.
  \item Ammonia-cracker reactor single-train cap 172 MMBtu/hr (distinct unit).
\end{itemize}

\begin{table}[htbp]
  \centering
  \caption{Per-unit single-train caps and replica counts. Number of parallel trains follows
    $n = \mathrm{ceil}(\text{total duty}/\text{cap})$; the cap equals the saturated per-replica load.}
  \label{tab:S3-1}
  \resizebox{\textwidth}{!}{%
  \begin{tabular}{@{}lllll@{}}
    \toprule
    Unit & Service & Sizing metric & Single-train cap & Replicas [1, 2.5, 5, 7.5, 10, 20, 30]$\times$ \\
    \midrule
    C-101 & Compressor & Volumetric flow & 15.81$\times$10\textsuperscript{6} cuft/hr (2772 HP) & 1 / 1 / 1 / 1 / 1 / 2 / 3 \\
    C-110 & Compressor & Volumetric flow & 17.57$\times$10\textsuperscript{6} cuft/hr & 1 / 1 / 1 / 2 / 2 / 3 / 5 \\
    C-147 & Compressor & Volumetric flow & 17.77$\times$10\textsuperscript{6} cuft/hr & 1 / 1 / 1 / 2 / 2 / 4 / 6 \\
    C-401 & Compressor & Volumetric flow & 15.23$\times$10\textsuperscript{6} cuft/hr (192,935 HP) & 1 / 1 / 1 / 1 / 1 / 2 / 3 \\
    MC-610 & Compressor & Volumetric flow & 17.38$\times$10\textsuperscript{6} cuft/hr & 1 / 2 / 3 / 4 / 5 / 10 / 20 $\dagger$ \\
    C-231 & Expander & Volumetric flow & 24,981 cuft/hr (1944 HP) & 2 / 6 / 10 / 14 / 18 / 36 / 54 $\dagger$ \\
    R-110 & Combustion chamber (pyrolyzer) & Heat absorbed & 447.5 MMBtu/hr (base 44.75) & 1 / 1 / 1 / 1 / 1 / 2 / 3 \\
    E-616 & Heat exchanger & Heat duty & 955.4 MMBtu/hr & 1 / 1 / 1 / 1 / 1 / 2 / 3 \\
    GT-612 & Gas turbine & Power output & 441,073 hp & 1 / 1 / 2 / 2 / 2 / 4 / 6 \\
    \bottomrule
  \end{tabular}}
\end{table}

\begin{itemize}
  \item The R-110 cap of 447.5 MMBtu/hr (10$\times$ the 44.75 MMBtu/hr base) supersedes the
    700 MMBtu/hr figure listed above; $\mathrm{ceil}(\text{scale} \times 44.75 / 447.5)$ reproduces its
    1/1/1/1/1/2/3 pattern exactly.
  \item $\dagger$ MC-610 and C-231 are \textbf{over-provisioned} relative to the strict
    $\mathrm{ceil}(\text{total}/\text{cap})$ rule: MC-610 carries 20 trains at 30$\times$ where the rule gives 15
    (total/cap = 15.0 exactly), and C-231 carries 6 and 10 trains at 2.5$\times$ and 5$\times$ where the rule
    gives 5 and 9. Every other unit matches the rule exactly.
\end{itemize}

%% ------------------------------------------------------------------
\section{Ammonia cracker detail}\label{sec:S4}

\begin{figure}[htbp]
    \centering
    \includegraphics[width=\textwidth]{BiomassLCA-Ammonia (2).png}
    \caption{Simplified process flow diagram of the ammonia cracking unit: green, blue, and orange line refer to the streams with ammonia feed, hydrogen-rich product, and residual gas, respectively.}
    \label{fig:pfd-ammonia}
\end{figure}

Ammonia cracker is simplified in Fig.~\ref{fig:pfd-ammonia}.
\begin{itemize}
  \item Cracker scale index $n$ = main \Htwo{} demand / 6570.40 lb/hr; non-grid $n$ linearly
    interpolated between bracketing APEA runs.
  \item Utilities = electricity; raw material = ammonia; fertilizer credited from the Product
    Summary sheet. Conversion / \Htwo-recovery ratios live inside the Aspen model (APEA export
    gives \NHthree-in vs \Htwo-out); GHSV/kinetics from \cite{cho2024ammonia}.
  \item Grey-\Htwo{} benchmark 0.68 \$/lb \Htwo. SMR uses natural-gas price 0.092306 \$/lb (US), escalated.
  \item 
\end{itemize}

%% ------------------------------------------------------------------
\section{CCS detail}\label{sec:S5}

\begin{itemize}
  \item Sequestered mass $\text{co2}_\text{seq} = \text{co2}_\text{left}(\text{623CO2})\cdot\text{ccs portion} + \text{co2}_\text{EGR}(\text{622X})$;
    622X captured at 100\%, 623CO2 at ccs portion.
  \item Capture OPEX: Fuel $= \max(\text{heat} - \text{plant heat gen}, 0)\cdot 14.38\cdot 8000$;
    Electricity $= \text{elec}\cdot 19.05\cdot 8000$; T\&S $= \text{co2}_\text{seq}/2204.62\cdot 8000\cdot(\text{co2 TS cost per ton}$).
  \item \COtwo{} T\&S price: default 0 (off) in base runs; literature range 4--45 \$/t; manuscript basis
    10 \$/t. 
  \item ccs portion sweeps: \{0, 0.25, 0.5, 0.75, 1.0\}; 0--0.9 in 0.1 steps; scale-up scenario fixed at 0.60.
\end{itemize}


%% ------------------------------------------------------------------
% \section{LCA inventories, allocation, and SAF criterion}\label{sec:S7}

% In this paper, LCA on the environmental impact of the developed process designs was conducted following the CORSIA methodology \cite{ICAO2023CORSIA} as illustrated in Fig. \ref{fig:lca-scope}.
% The International Civil Aviation Organization (ICAO) has proposed the method as a standard so that the SAF industry can consistently and transparently monitor its GHG emission intensities. 
% % It was formerly proposed by the International Civil Aviation Organization (ICAO) as a standard to guide the SAF industry to monitor their GHG emission intensities consistently and transparently. 
% CORSIA life cycle emission value is the sum of the direct emissions along the supply chains of individual SAFs and the estimated emissions for induced land use change (ILUC) \cite{prussi2021corsia}.
% It is also referred to as WTWa emissions since its scope spans from feedstock preparation to fuel combustion on the aircraft. 



%% ------------------------------------------------------------------
\section{Biomass price model in supply-availability context}\label{sec:S7}

Domestic feedstock ceilings used to convert plant scale to a share of aviation-fuel demand: forest residue 3,127,378 t/yr; log (2,843,071 + 3,274,200) t/yr. 

\textbf{FR Price (reference).} A constant forest-residue price of 107.3~\$/t, independent of scale; this is the optimistic floor.

\textbf{FR--Log Linear} (naive Galik supply curve; \citealp{galik}). The forest-residue price holds until regional forest-residue supply (2.84~MMt/yr) is exhausted, beyond which the delivered price rises linearly with demand, price = supply $\times$ 25.08 + 33.30~\$/t (supply in MMt/yr), as logging residue is drawn in.

\textbf{Fixed Elasticity} (Galik supply curve, supply elasticity 0.3). Follows the FR--Log Linear curve until the combined forest-residue-plus-log supply (6.12~MMt/yr) is exhausted, then extrapolates linearly under a fixed supply elasticity of 0.3 (log reference price 192.4~\$/t), giving the steepest scale penalty.

\textbf{Scale-up evaluation.} MJSP is evaluated at seven plant scales: 0.041, 0.104, 0.207, 0.311, 0.414, 0.828 and 1.243~MMt/yr SAF. At the forest-residue floor price, the base-case (grey-\ce{H2}) MJSP at each scale is broken down into capital investment, raw materials, utilities, fixed operating costs (operating labor, maintenance, operating charges, plant overhead, and general and administrative costs) and co-product sales, which enter as a negative term (Table~\ref{tab:S9-1}). Each of the three supply-curve models then sets the delivered biomass price at each scale, and MJSP is recomputed for three configurations: the grey-\ce{H2} base, full green ammonia, and full CCS (Tables~\ref{tab:S9-2} and \ref{tab:S9-3}). Biomass is costed identically in all three configurations, so a given price model shifts each of them by the same biomass-cost increment. The minimum-MJSP scale is read where the falling conversion cost and the rising feedstock cost cross.
%% ------------------------------------------------------------------
%% ------------------------------------------------------------------
\section{Data tables for the Results figures}\label{sec:S9}

Numeric values behind the main-text Results figures.
MJSP in \$/L SAF (1Q-2022 USD); GHG in g \COtwoe/MJ (WTWa GHG-100, Allocation, Korea grid mix).

\begin{table}[htbp]
  \centering
  \caption{Base-case (grey \Htwo, flat forest-residue price) MJSP breakdown vs SAF production
    (base-case breakdown figure).
    Components in \$/L SAF; MJSP = sum of components.}
  \label{tab:S9-1}
  \resizebox{\textwidth}{!}{%
  \begin{tabular}{@{}rrrrrrrrrrr@{}}
    \toprule
    \shortstack{SAF production\\{[MMt/yr]}} & \shortstack{Capital\\investment} & \shortstack{Raw\\materials} & Utilities &
    \shortstack{Operating\\labor} & Maintenance & \shortstack{Operating\\charges} & \shortstack{Plant\\overhead} & G\&A &
    \shortstack{Co-product\\sales} & MJSP \\
    \midrule
    0.041 & 1.152 & 1.284 & 0.676 & 0.036 & 0.110 & 0.009 & 0.073 & 0.172 & $-1.860$ & 1.651 \\
    0.104 & 0.936 & 1.269 & 0.603 & 0.014 & 0.093 & 0.004 & 0.054 & 0.160 & $-1.860$ & 1.271 \\
    0.207 & 0.857 & 1.263 & 0.480 & 0.007 & 0.086 & 0.002 & 0.046 & 0.147 & $-1.860$ & 1.028 \\
    0.311 & 0.801 & 1.262 & 0.419 & 0.005 & 0.083 & 0.001 & 0.044 & 0.142 & $-1.860$ & 0.897 \\
    0.414 & 0.750 & 1.261 & 0.381 & 0.004 & 0.078 & 0.001 & 0.041 & 0.138 & $-1.860$ & 0.793 \\
    0.828 & 0.719 & 1.260 & 0.340 & 0.003 & 0.077 & 0.001 & 0.040 & 0.134 & $-1.860$ & 0.714 \\
    1.243 & 0.700 & 1.259 & 0.324 & 0.002 & 0.079 & 0.000 & 0.040 & 0.133 & $-1.860$ & 0.676 \\
    \bottomrule
  \end{tabular}}
\end{table}

\begin{table}[htbp]
  \centering
  \caption{Biomass price and grey-\Htwo{} base MJSP under three biomass price models
    (price-model figure). MJSP in \$/L SAF.}
  \label{tab:S9-2}
  \resizebox{\textwidth}{!}{%
  \begin{tabular}{@{}rrrrrrr@{}}
    \toprule
    & \multicolumn{3}{c}{Biomass price [\$/lb]} & \multicolumn{3}{c}{MJSP grey \Htwo{} [\$/L]} \\
    \cmidrule(lr){2-4}\cmidrule(l){5-7}
    SAF production [MMt/yr] & FR Price & FR-Log Linear & Fixed Elasticity & FR Price & FR-Log Linear & Fixed Elasticity \\
    \midrule
    0.041 & 0.0486 & 0.0486 & 0.0486 & 1.651 & 1.651 & 1.651 \\
    0.104 & 0.0486 & 0.0486 & 0.0486 & 1.271 & 1.271 & 1.271 \\
    0.207 & 0.0486 & 0.0486 & 0.0486 & 1.028 & 1.028 & 1.028 \\
    0.311 & 0.0486 & 0.0486 & 0.0486 & 0.897 & 0.897 & 0.897 \\
    0.414 & 0.0486 & 0.0527 & 0.0527 & 0.793 & 0.849 & 0.849 \\
    0.828 & 0.0486 & 0.0903 & 0.1064 & 0.714 & 1.290 & 1.515 \\
    1.243 & 0.0486 & 0.1278 & 0.2360 & 0.676 & 1.773 & 3.272 \\
    \bottomrule
  \end{tabular}}
\end{table}

\begin{table}[htbp]
  \centering
  \caption{MJSP of the full green-\NHthree{} and full-CCS configurations under three biomass price models
    (price-model figure; full green \NHthree{} = 100\% green-\NHthree{} hydrogen,
    full CCS = CCS portion 1.0). MJSP in \$/L SAF.}
  \label{tab:S9-3}
  \resizebox{\textwidth}{!}{%
  \begin{tabular}{@{}rrrrrrr@{}}
    \toprule
    & \multicolumn{3}{c}{MJSP full green \NHthree{} [\$/L]} & \multicolumn{3}{c}{MJSP full CCS [\$/L]} \\
    \cmidrule(lr){2-4}\cmidrule(l){5-7}
    SAF production [MMt/yr] & FR Price & FR-Log Linear & Fixed Elasticity & FR Price & FR-Log Linear & Fixed Elasticity \\
    \midrule
    0.041 & 3.495 & 3.495 & 3.495 & 2.275 & 2.275 & 2.275 \\
    0.104 & 3.062 & 3.062 & 3.062 & 1.852 & 1.852 & 1.852 \\
    0.207 & 2.802 & 2.802 & 2.802 & 1.586 & 1.586 & 1.586 \\
    0.311 & 2.661 & 2.661 & 2.661 & 1.445 & 1.445 & 1.445 \\
    0.414 & 2.544 & 2.600 & 2.600 & 1.334 & 1.390 & 1.390 \\
    0.828 & 2.443 & 3.020 & 3.244 & 1.242 & 1.818 & 2.043 \\
    1.243 & 2.401 & 3.498 & 4.997 & 1.198 & 2.295 & 3.794 \\
    \bottomrule
  \end{tabular}}
\end{table}

\begin{table}[htbp]
  \centering
  \caption{MJSP, GHG, Pareto status, sensitivity band, and abatement cost of the 30 green-\NHthree{} share
    $\times$ CCS portion configurations at 0.414 MMt/yr SAF.  Abatement cost relative to CCS 60\%
    (75.14 g \COtwoe/MJ, \$0.977/L); -- where GHG is not below CCS 60\%.}
  \label{tab:S9-4}
  \small
  \begin{tabular}{@{}rrrrcr@{}}
    \toprule
    \shortstack{Green \NHthree\\share [\%]} & \shortstack{CCS\\portion [\%]} & \shortstack{MJSP\\{[\$/L]}} &
    \shortstack{GHG\\{[g \COtwoe/MJ]}} & Pareto &
    \shortstack{Abatement cost\\{[\$/t \COtwo]}} \\
    \midrule
    0 & 0 & 0.792 & 100.83 & yes & -- \\
    0 & 25 & 0.855 & 89.93 & yes & -- \\
    0 & 50 & 0.913 & 78.99 & yes & -- \\
    0 & 75 & 1.111 & 69.36 & yes & 733 \\
    0 & 100 & 1.334 & 59.96 & yes & 742 \\
    10 & 0 & 0.977 & 97.02 & no & -- \\
    10 & 25 & 1.039 & 86.12 & no & -- \\
    10 & 50 & 1.098 & 75.18 & yes & -- \\
    10 & 75 & 1.296 & 65.55 & yes & 1,049 \\
    10 & 100 & 1.518 & 56.15 & yes & 900 \\
    25 & 0 & 1.240 & 91.31 & no & -- \\
    25 & 25 & 1.302 & 80.40 & no & -- \\
    25 & 50 & 1.361 & 69.46 & no & 2,138 \\
    25 & 75 & 1.559 & 59.83 & no & 1,201 \\
    25 & 100 & 1.782 & 50.44 & yes & 1,029 \\
    50 & 0 & 1.680 & 81.78 & no & -- \\
    50 & 25 & 1.742 & 70.88 & no & 5,671 \\
    50 & 50 & 1.800 & 59.94 & no & 1,711 \\
    50 & 75 & 1.998 & 50.31 & yes & 1,299 \\
    50 & 100 & 2.221 & 40.91 & yes & 1,148 \\
    75 & 0 & 2.116 & 72.26 & no & 12,511 \\
    75 & 25 & 2.178 & 61.35 & no & 2,752 \\
    75 & 50 & 2.236 & 50.41 & no & 1,609 \\
    75 & 75 & 2.434 & 40.78 & yes & 1,340 \\
    75 & 100 & 2.657 & 31.39 & yes & 1,213 \\
    100 & 0 & 2.544 & 62.74 & no & 3,991 \\
    100 & 25 & 2.606 & 51.83 & no & 2,207 \\
    100 & 50 & 2.665 & 40.89 & no & 1,556 \\
    100 & 75 & 2.862 & 31.26 & yes & 1,357 \\
    100 & 100 & 3.085 & 21.86 & yes & 1,250 \\
    \bottomrule
  \end{tabular}
\end{table}

\begin{table}[htbp]
  \centering
  \caption{Grey-\Htwo{} price at which Green \NHthree{} 10\% + CCS 50\% matches CCS 60\%.}
  \label{tab:S9-5}
  \small
  \begin{tabular}{@{}rrrc@{}}
    \toprule
    Green \NHthree{} price [\$/t] & Crossover grey-\Htwo{} price [\$/kg] & MJSP at crossover [\$/L] & Inside 1--4 \$/kg range \\
    \midrule
    270 & 1.59 & 1.016 & yes \\
    441 & 2.91 & 1.568 & yes \\
    635 (base) & 4.40 & 2.195 & no \\
    1,000 & 7.21 & 3.373 & no \\
    \bottomrule
  \end{tabular}
\end{table}

\begin{table}[htbp]
  \centering
  \caption{Mean $|$SHAP$|$ of MJSP drivers at the two cheapest compliant points in \$/L SAF: 128 Latin-hypercube TEA runs per point,
    gradient-boosted surrogate ($R^2 > 0.999$); -- = driver not sampled.}
  \label{tab:S9-6}
  \small
  \begin{tabular}{@{}lrr@{}}
    \toprule
    Driver & CCS 60\% & Green \NHthree{} 10\% + CCS 50\% \\
    \midrule
    Grey \Htwo & 0.310 & 0.274 \\
    CAPEX & 0.103 & 0.102 \\
    Co-product & 0.097 & 0.083 \\
    Electricity & 0.052 & 0.055 \\
    Green \NHthree & -- & 0.031 \\
    \COtwo{} T\&S & 0.025 & 0.023 \\
    Capacity & 0.019 & 0.028 \\
    Biomass & 0.003 & 0.002 \\
    \bottomrule
  \end{tabular}
\end{table}

\begin{table}[htbp]
  \centering
  \caption{One-at-a-time MJSP spans at the two cheapest compliant points in \$/L SAF; input ranges as in the tornado-driver analysis.}
  \label{tab:S9-7}
  \small
  \begin{tabular}{@{}llrrrr@{}}
    \toprule
    Point & Driver & Base MJSP & MJSP at low input & MJSP at high input & Span \\
    \midrule
    CCS 60\% & Grey \Htwo & 0.977 & 0.767 & 2.026 & 1.258 \\
    CCS 60\% & Co-product & 0.977 & 1.224 & 0.730 & 0.494 \\
    CCS 60\% & CAPEX & 0.977 & 0.741 & 1.213 & 0.472 \\
    CCS 60\% & Electricity & 0.977 & 1.134 & 0.820 & 0.313 \\
    CCS 60\% & Capacity & 0.977 & 1.174 & 0.987 & 0.188 \\
    CCS 60\% & \COtwo{} T\&S & 0.977 & 0.950 & 1.137 & 0.187 \\
    CCS 60\% & Biomass & 0.977 & 0.977 & 1.029 & 0.052 \\
    \midrule
    Green \NHthree{} 10\% + CCS 50\% & Grey \Htwo & 1.098 & 0.909 & 2.042 & 1.133 \\
    Green \NHthree{} 10\% + CCS 50\% & Co-product & 1.098 & 1.344 & 0.851 & 0.494 \\
    Green \NHthree{} 10\% + CCS 50\% & CAPEX & 1.098 & 0.858 & 1.337 & 0.480 \\
    Green \NHthree{} 10\% + CCS 50\% & Electricity & 1.098 & 1.258 & 0.937 & 0.321 \\
    Green \NHthree{} 10\% + CCS 50\% & Green \NHthree & 1.098 & 0.980 & 1.215 & 0.234 \\
    Green \NHthree{} 10\% + CCS 50\% & Capacity & 1.098 & 1.313 & 1.108 & 0.205 \\
    Green \NHthree{} 10\% + CCS 50\% & \COtwo{} T\&S & 1.098 & 1.073 & 1.238 & 0.165 \\
    Green \NHthree{} 10\% + CCS 50\% & Biomass & 1.098 & 1.149 & 1.149 & 0.052 \\
    \bottomrule
  \end{tabular}
\end{table}

\bibliographystyle{cas-model2-names}
\bibliography{main-sc}

\end{document}

